\section{课程导读：从 AI 工具到 AI 共研者}

这节课的核心不只是``AI 可以做科研''，而是科研流程的组织方式正在变化。James Zou 给出的叙事是三段式：第一段是 Virtual Lab（多 Agent 科学团队），第二段是 Paper2Agent（把论文转为可交互 Agent），第三段是 Agents for Science（用真实会议数据观察人机协作与 AI 评审的行为）。

如果把过去十年的科学 AI 记为``模型能力持续增强''，那么这一讲强调的是另一个维度：\textbf{科学工作流的可编排性}。在这个维度里，问题不再是``某个单模型能不能刷高 benchmark''，而是``能不能把一个研究课题从问题定义、知识检索、方案生成、计算验证到结果解释，组织成一个可迭代的 agentic pipeline''。

\begin{importantbox}{本讲主线}
\begin{itemize}
    \item 从``单任务 AI 工具''转向``端到端研究协作系统''。
    \item 核心对象从模型参数变成\textbf{角色分工 + 交互协议 + 外部工具接入}。
    \item 评估指标从单点准确率扩展到\textbf{可复现性、错误恢复能力、人与 Agent 的分工质量}。
\end{itemize}
\end{importantbox}

Zou 在开场里反复强调一个概念转移：过去我们通常先有定义明确的问题，再为该问题选一个专用模型；而现在越来越多团队尝试让 Agent 参与更上游、更开放的研究环节。这个变化对应一句非常关键的话：\textit{``people are starting to explore AI as a co-scientist''}。这句话直接把 Agent 的角色从``执行器''抬到了``协作者''。

\begin{knowledgebox}{为何 2025 年后这个转移突然加速}
\begin{itemize}
    \item LLM 的工具调用、长上下文、代码执行能力形成了基础设施闭环。
    \item MCP 等协议让``资源如何被模型稳定调用''有了工程标准化路径。
    \item 科学研究本身是高复杂度、多阶段任务，天然适合多 Agent 协作分治。
\end{itemize}
\end{knowledgebox}

\begin{warningbox}{常见误解：把 Agent 当成更大参数模型}
这节课展示的系统优势大多不来自``更大模型''，而来自编排层设计：角色定义、会议机制、记忆机制、sandbox、critic 角色、human-in-the-loop 纠偏。忽略这些机制，只比较底座模型，往往会得出错误结论。
\end{warningbox}

\subsection{本章小结}

本章建立了整节课的理解框架：这不是单点模型能力展示，而是科研系统工程范式。后续章节可按``组织机制 -> 案例验证 -> 平台化扩展 -> 会议级证据 -> 局限与治理''来阅读。

